\documentclass[12pt]{article}

\usepackage[spanish]{babel}
\usepackage[utf8]{inputenc}
\usepackage[T1]{fontenc}
\usepackage{geometry}
\usepackage{graphicx}
\usepackage{setspace}
\usepackage{titlesec}
\usepackage{listings}
\usepackage{xcolor}
\usepackage{float}
\usepackage{hyperref}

\geometry{margin=1in}
\onehalfspacing

\definecolor{codegray}{rgb}{0.95,0.95,0.95}

\lstset{
    backgroundcolor=\color{codegray},
    basicstyle=\ttfamily\small,
    breaklines=true,
    frame=single,
    language=Python
}

\begin{document}

% portada

\begin{titlepage}
    \centering

    %\includegraphics[width=0.3\textwidth]{logo.png}\par\vspace{1cm}

    {\scshape\Large Bases de Datos NoSQL \par}
    \vspace{1.5cm}

    {\huge\bfseries Evidencia de la Práctica CRUD \par}
    \vspace{0.5cm}

    {\Large Aplicación Web con Flask y MongoDB Atlas \par}

    \vfill

    {\large
    Nombre del Alumno: Emilio Izquierdo Montero \\
    No Control: 22660231 \\
    Nombre del Alumno: Acosta Cardenas Eduardo Santiago \\
    No Control: 22660221 \\
    \vspace{0.3cm}

    Profesor: Carlos Alberto Garcia Sanchez \\
    \vspace{0.3cm}

    Grupo / Semestre: 8 \\
    \vspace{0.3cm}

    Fecha: 8 de mayo del 2026
    \par}

    \vfill

\end{titlepage}

\section{Introducción}

La presente práctica consistió en el desarrollo de una aplicación web funcional utilizando Python, Flask y MongoDB Atlas como base de datos NoSQL en la nube.

La aplicación implementa las operaciones CRUD (Create, Read, Update y Delete) sobre una colección de restaurantes obtenida desde un dataset real en formato JSON. El sistema permite visualizar restaurantes, realizar búsquedas, registrar nuevos establecimientos, actualizar información y eliminar registros desde una interfaz web.

El proyecto también implementa conexión remota a MongoDB Atlas utilizando PyMongo y variables de entorno para proteger la URI de conexión.


\section{Objetivo}

Desarrollar una aplicación web conectada a MongoDB Atlas capaz de:

\begin{itemize}
    \item Visualizar restaurantes almacenados en MongoDB.
    \item Filtrar restaurantes por distrito o tipo de cocina.
    \item Registrar nuevos restaurantes.
    \item Editar restaurantes existentes.
    \item Eliminar registros mediante confirmación.
\end{itemize}


\section{Tecnologías Utilizadas}

Las tecnologías y herramientas utilizadas fueron:

\begin{itemize}
    \item Python 3
    \item Flask
    \item PyMongo
    \item MongoDB Atlas
    \item HTML
    \item CSS
    \item dotenv
\end{itemize}


\section{Librerías Utilizadas}

La aplicación utiliza las siguientes librerías principales:

\begin{lstlisting}
from flask import Flask, render_template,
request, redirect, url_for, jsonify, flash

from pymongo import MongoClient

from bson import ObjectId

from dotenv import load_dotenv
\end{lstlisting}

\begin{itemize}
    \item Flask: Framework web utilizado para la interfaz y rutas.
    \item PyMongo: Librería para conectarse y manipular MongoDB.
    \item bson: Permite trabajar con ObjectId.
    \item dotenv: Carga variables de entorno desde un archivo .env.
\end{itemize}


\section{Conexión con MongoDB Atlas}

La conexión hacia MongoDB Atlas se realizó utilizando una URI almacenada dentro de variables de entorno.

\begin{lstlisting}
load_dotenv()

MONGO_URI = os.getenv("MONGO_URI")

client = MongoClient(MONGO_URI)

db = client["restaurantDB"]

collection = db["restaurants"]
\end{lstlisting}

El uso de variables de entorno permite proteger credenciales y mantener una configuración más segura.


\section{Carga del Dataset}

Se desarrolló un script independiente llamado \texttt{import\_data.py} encargado de importar el dataset de restaurantes hacia MongoDB Atlas.

El script:

\begin{itemize}
    \item Lee el archivo JSON.
    \item Convierte cada línea en documentos JSON.
    \item Inserta los datos en lotes.
    \item Permite limpiar la colección antes de importar nuevamente.
\end{itemize}

\begin{lstlisting}
records = load_jsonl(filepath)

client = MongoClient(MONGO_URI)

db = client[DB_NAME]

col = db[COL_NAME]

col.insert_many(batch)
\end{lstlisting}


\section{Estructura del Dataset}

Cada documento del dataset contiene información relacionada con restaurantes.

Ejemplo de documento:

\begin{lstlisting}
{
    "name": "Pizza Town",
    "borough": "Brooklyn",
    "cuisine": "Pizza",
    "address": {
        "building": "100",
        "street": "Main Street",
        "zipcode": "11201"
    }
}
\end{lstlisting}


\section{Operaciones CRUD Implementadas}

\subsection{READ — Visualización y Búsqueda}

La aplicación permite listar restaurantes mostrando:

\begin{itemize}
    \item Nombre
    \item Distrito
    \item Tipo de cocina
\end{itemize}

Además, se implementaron filtros de búsqueda utilizando parámetros GET.

\begin{lstlisting}
borough = request.args.get("borough", "")
cuisine = request.args.get("cuisine", "")

query = {}

if borough:
    query["borough"] = borough

if cuisine:
    query["cuisine"] = cuisine

restaurants = [
    serialize(r)
    for r in collection.find(query).limit(100)
]
\end{lstlisting}


\subsection{CREATE — Registro de Restaurantes}

La aplicación incluye un formulario para registrar nuevos restaurantes.

Campos utilizados:

\begin{itemize}
    \item Nombre
    \item Distrito
    \item Tipo de cocina
    \item Código Postal
    \item Calle
    \item Número
\end{itemize}

\begin{lstlisting}
doc = {
    "name": name,
    "borough": borough,
    "cuisine": cuisine,
    "address": {
        "building": building,
        "street": street,
        "zipcode": zipcode,
        "coord": [],
    },
    "grades": [],
}

collection.insert_one(doc)
\end{lstlisting}


\subsection{UPDATE — Actualización de Datos}

La aplicación permite modificar restaurantes mediante su identificador.

Campos actualizables:

\begin{itemize}
    \item Nombre
    \item Tipo de cocina
    \item Distrito
    \item Código Postal
\end{itemize}

\begin{lstlisting}
collection.update_one(
    {"_id": oid},
    {
        "set": update_fields
    }
)
\end{lstlisting}


\subsection{DELETE — Eliminación de Restaurantes}

Se implementó una función para eliminar restaurantes mediante confirmación.

\begin{lstlisting}
result = collection.delete_one({
    "_id": oid
})
\end{lstlisting}


\section{API JSON}

La aplicación también incluye una pequeña API REST para devolver restaurantes en formato JSON.

\begin{lstlisting}
@app.route("/api/restaurants")
def api_list():

    docs = [
        serialize(d)
        for d in collection.find(query).limit(50)
    ]

    return jsonify(docs)
\end{lstlisting}

Esta API fue utilizada principalmente para pruebas y depuración.


\section{Evidencias de Ejecución}

\subsection{Frontend de la aplicación}

Captura donde se observa el Frontend de la aplicación
\begin{figure}[H]
    \centering
    \includegraphics[width=0.95\textwidth]{/home/emilio/Descargas/frontendrestaurante.png}
\end{figure}


\subsection{Agregar un restaurante}
Restaurante siendo insertado con el endpoint en la base de datos 

\begin{figure}[H]
    \centering
    \includegraphics[width=0.95\textwidth]{/home/emilio/Descargas/insertar.png}
\end{figure}


\subsection{Buscado del restaurante insertado}

Visualización y búsqueda de restaurantes.

\begin{figure}[H]
    \centering
    \includegraphics[width=0.95\textwidth]{/home/emilio/Descargas/busqueda.png}
\end{figure}



\subsection{Edición de Restaurante}

Actualización de datos desde la interfaz web.

\begin{figure}[H]
    \centering
    \includegraphics[width=0.95\textwidth]{/home/emilio/Descargas/editar.png}
\end{figure}


\subsection{Eliminación de Restaurante}

Eliminación de registros mediante confirmación.

\begin{figure}[H]
    \centering
    \includegraphics[width=0.95\textwidth]{/home/emilio/Descargas/borrar1.png}
\end{figure}

\begin{figure}[H]
    \centering
    \includegraphics[width=0.95\textwidth]{/home/emilio/Descargas/borrar2.png}
\end{figure}


\section{Archivo Requirements}

El proyecto utiliza las siguientes dependencias:

\begin{lstlisting}
flask
pymongo
python-dotenv
dnspython
\end{lstlisting}


\section*{Código Fuente}

El código fuente del proyecto se encuentra disponible en GitHub:

\begin{itemize}
    \item Repositorio:
    \href{https://github.com/emili0p/CocinasManhattan}
    {https://github.com/emili0p/CocinasManhattan}
\end{itemize}

\noindent
El repositorio incluye:

\begin{itemize}
    \item Código fuente completo
    \item Templates HTML
    \item Archivo requirements.txt
    \item Script import\_data.py
\end{itemize}


\section{Conclusión}

La práctica permitió comprender el funcionamiento de bases de datos NoSQL utilizando MongoDB Atlas y su integración con aplicaciones web desarrolladas con Flask.

El uso de PyMongo facilitó las operaciones CRUD sobre documentos JSON, mientras que Flask permitió construir una interfaz sencilla y funcional para administrar la información de restaurantes.

Además, el uso de variables de entorno mediante dotenv ayudó a mejorar la seguridad de la conexión hacia MongoDB Atlas.

\end{document}
